\documentclass[conference]{IEEEtran}
\IEEEoverridecommandlockouts

\usepackage{cite}
\usepackage{amsmath,amssymb,amsfonts}
\usepackage{algorithmic}
\usepackage{graphicx}
\usepackage{textcomp}
\usepackage{xcolor}
\usepackage{booktabs}
\usepackage{array}
\usepackage{multirow}
\usepackage{url}

\def\BibTeX{{\rm B\kern-.05em{\sc i\kern-.025em b}\kern-.08em
    T\kern-.1667em\lower.7ex\hbox{E}\kern-.125emX}}

\begin{document}

\title{BlockShield: A Hybrid Blockchain and Explainable Machine Learning
Framework for Certificate Fraud Detection in Online Learning Platforms}

\author{
\IEEEauthorblockN{[Your Name]}
\IEEEauthorblockA{\textit{PG Scholar, Dept. of Computer Science \& Engineering} \\
\textit{[Your College Name]} \\
\mbox{[City]}, India \\
\texttt{your.email@college.edu}}
\and
\IEEEauthorblockN{[Guide Name]}
\IEEEauthorblockA{\textit{Assistant Professor, Dept. of Computer Science \& Engineering} \\
\textit{[Your College Name]} \\
\mbox{[City]}, India \\
\texttt{guide.email@college.edu}}
}

\maketitle

\begin{abstract}
The rapid expansion of online learning has made course completion
certificates a widely accepted credential in hiring and academic
admission processes. This growing reliance has simultaneously
attracted fraudulent activity, including certificate forgery,
unauthorised reuse, and coordinated bot-driven verification attacks.
Conventional certificate management systems depend on centralised
storage, which introduces a single point of failure and offers no
mechanism to detect post-issuance misuse. This paper introduces
\textit{BlockShield}, a two-tier fraud mitigation architecture that
anchors each certificate's cryptographic identity on the Ethereum
Sepolia blockchain while concurrently applying an XGBoost-based
behavioural fraud classifier to score every certificate in real time.
The classifier operates on 24 engineered features spanning completion
behaviour, verification patterns, and blockchain metadata. SHapley
Additive exPlanations (SHAP) decompose each prediction into
per-feature contributions, enabling administrators to understand
precisely which signals triggered a fraud alert. A GraphSAGE Graph
Neural Network is trained on the same dataset for comparative
evaluation. On a 1,000-sample domain-informed synthetic dataset,
XGBoost attains AUC\,=\,1.00 and F1\,=\,1.00, outperforming
GraphSAGE (AUC\,=\,0.997, F1\,=\,0.967). The complete system is
deployed as a production web application on cloud infrastructure,
incorporating a circuit-breaker fault-tolerance mechanism and an
automated model retraining pipeline.
\end{abstract}

\begin{IEEEkeywords}
certificate fraud detection, blockchain, XGBoost, SHAP, explainable
artificial intelligence, GraphSAGE, Ethereum, smart contracts,
online learning, credential verification
\end{IEEEkeywords}

% ─────────────────────────────────────────────────────────────────
\section{Introduction}

\subsection{Problem Context}

Online education platforms now issue millions of course completion
certificates annually. Employers and academic institutions
increasingly treat these credentials as evidence of skill
acquisition, creating strong incentives for fraudulent actors to
fabricate, alter, or misuse them. Three categories of abuse are
particularly prevalent: (i) direct document forgery, where a
certificate image is edited to change the recipient name or course
title; (ii) database tampering, where an attacker with access to
the issuing platform modifies stored records; and (iii) behavioural
fraud, where a certificate is obtained through automated means such
as bot-driven quiz completion or credential stuffing attacks during
the verification phase.

Existing countermeasures address these threats in isolation.
QR-code-based verification links back to a centralised server,
which itself can be compromised. Watermarking deters image editing
but cannot detect database-level manipulation. Neither approach
provides any signal about suspicious verification behaviour after
issuance.

\subsection{Proposed Approach}

BlockShield addresses all three fraud categories within a single
integrated architecture. Blockchain registration eliminates
post-issuance tampering by anchoring an immutable SHA-256 hash of
each certificate on the Ethereum Sepolia testnet at the moment of
issuance. Any subsequent modification to the certificate data
produces a different hash, making tampering cryptographically
detectable without relying on any centralised authority.

Behavioural fraud, which blockchain alone cannot detect, is handled
by a machine learning layer. An XGBoost classifier continuously
monitors 24 features extracted from certificate metadata and
verification event logs. Each prediction is accompanied by SHAP
values that identify the top contributing features, giving
administrators an interpretable audit trail rather than an opaque
score.

\subsection{Research Questions}

This work investigates three questions:
\begin{enumerate}
    \item Does on-chain hash registration provide reliable tamper
    detection in a production deployment?
    \item Which model — XGBoost or GraphSAGE GNN — better classifies
    certificate fraud on structured tabular verification data?
    \item Can SHAP values produce actionable, per-certificate
    explanations suitable for non-technical administrative review?
\end{enumerate}

\subsection{Contributions}

The key contributions of this paper are:
\begin{itemize}
    \item A production-deployed system that unifies Ethereum smart
    contracts with real-time ML fraud scoring in a single platform.
    \item A 24-feature behavioural fraud model with SHAP-driven
    explainability surfaced through an admin monitoring dashboard.
    \item An empirical comparison of XGBoost and GraphSAGE GNN on
    certificate verification data, with analysis of why tabular
    models outperform graph models at this scale.
    \item A retraining pipeline and circuit-breaker mechanism that
    maintain system reliability as real data accumulates.
\end{itemize}

% ─────────────────────────────────────────────────────────────────
\section{Related Work}

\subsection{Blockchain-Based Credential Systems}

The application of distributed ledger technology to academic
credentialing has been explored since the mid-2010s. The Blockcerts
initiative~\cite{ref1} defined an open standard for issuing
verifiable certificates on Bitcoin and Ethereum, focusing on
issuance integrity and recipient-controlled sharing. Grech and
Camilleri~\cite{ref2} surveyed blockchain applications in education
and identified immutability and disintermediation as the primary
advantages over traditional registries. A key limitation of these
systems is that they treat verification as a binary pass/fail
operation and provide no mechanism to detect anomalous verification
behaviour after a certificate has been legitimately issued.

\subsection{Machine Learning for Fraud Classification}

Gradient boosted decision trees have consistently achieved
state-of-the-art results on structured fraud detection tasks.
Chen and Guestrin~\cite{ref3} introduced XGBoost and demonstrated
its efficiency advantages over conventional boosting methods on
large-scale classification benchmarks. Subsequent studies in
financial fraud detection~\cite{ref4} confirmed that XGBoost
handles class imbalance and high-dimensional tabular features more
effectively than deep neural networks when training data is limited.

Graph-based fraud detection has gained attention as a complementary
approach. Hamilton et al.~\cite{ref5} proposed GraphSAGE, an
inductive graph representation learning method that generates node
embeddings by aggregating features from sampled neighbourhoods.
This makes it suitable for detecting coordinated fraud rings where
individual records appear legitimate but exhibit suspicious
relational patterns when viewed collectively.

\subsection{Explainable AI in Security Applications}

Lundberg and Lee~\cite{ref6} introduced SHAP as a theoretically
grounded framework for model explanation, deriving feature
attributions from Shapley values in cooperative game theory. Unlike
earlier importance measures such as permutation importance or
gradient-based saliency, SHAP guarantees consistency and local
accuracy. Applications in cybersecurity include network intrusion
detection~\cite{ref7} and credit risk assessment~\cite{ref8}, where
SHAP explanations have been shown to improve analyst trust and
reduce false-positive investigation time.

\subsection{Gap Addressed}

No prior work combines on-chain certificate registration with
continuous ML-based behavioural monitoring and SHAP explainability
in a single deployed system. BlockShield fills this gap by treating
blockchain and ML as complementary layers rather than alternative
solutions.

% ─────────────────────────────────────────────────────────────────
\section{System Architecture}

\subsection{Architectural Overview}

BlockShield is structured as three cooperating layers, illustrated
conceptually in Fig.~\ref{fig:architecture}:

\begin{enumerate}
    \item \textbf{Blockchain Layer} — Ethereum Sepolia smart contract
    for immutable certificate hash registration.
    \item \textbf{AI Scoring Layer} — FastAPI microservice hosting
    the XGBoost model and SHAP explainer.
    \item \textbf{Application Layer} — Node.js/Express REST API,
    React web frontend, and admin monitoring dashboard.
\end{enumerate}

\begin{figure}[htbp]
\centerline{\includegraphics[width=\columnwidth]{architecture.png}}
\caption{BlockShield three-layer architecture. The application layer
handles user interactions; the blockchain layer provides tamper
detection; the AI layer provides behavioural fraud scoring.}
\label{fig:architecture}
\end{figure}

\subsection{Certificate Issuance Flow}

A certificate is issued only after a user satisfies three
conditions: all course lessons are marked complete, all mandatory
post-lesson quizzes are passed, and no prior certificate exists for
the same user-course pair. The eligibility check queries
\texttt{CourseProgress} records and \texttt{MicroQuizAttempt}
records, supporting both the current modular lesson structure and
a legacy embedded-lesson format.

Upon eligibility confirmation, the system:
\begin{enumerate}
    \item Creates an immutable \texttt{Certificate} document in
    MongoDB, computing a SHA-256 hash over the tuple
    $\langle$\textit{certificateId, userId, courseId, issuedAt}$\rangle$.
    \item Submits a blockchain transaction calling
    \texttt{registerCertificate(id, hash)} on the deployed
    \texttt{CertificateRegistry} Solidity contract.
    \item Generates a PDF certificate using a Puppeteer-based
    renderer with a simple-PDF fallback.
    \item Triggers an initial fraud score computation via the AI
    microservice.
\end{enumerate}

\subsection{Blockchain Layer}

The \texttt{CertificateRegistry} contract is deployed on Ethereum
Sepolia. Its storage maps each certificate identifier to a
\texttt{Certificate} struct containing the 32-byte hash and the
block timestamp. Two invariants are enforced at the contract level:
a certificate identifier may be registered only once, and a given
hash value may appear only once across all registrations. These
constraints prevent both re-registration attacks and hash collision
exploits.

Verification calls \texttt{verifyCertificate(id, hash)}, a
\texttt{view} function that costs no gas. A mismatch between the
stored hash and the recomputed hash of the current certificate data
constitutes cryptographic proof of tampering, independent of the
application database.

\subsection{Feature Engineering}

The fraud classifier operates on 24 numerical features grouped into
four categories, as summarised in Table~\ref{tab:features}.

\begin{table}[htbp]
\caption{Feature Groups Used by the Fraud Classifier}
\label{tab:features}
\begin{center}
\begin{tabular}{lcp{4.2cm}}
\toprule
\textbf{Group} & \textbf{Count} & \textbf{Representative Features} \\
\midrule
Certificate Metadata   & 6 & revocation flag, blockchain tx presence,
                             wallet reuse count, grade encoding \\
Completion Behaviour   & 7 & completion \%, lesson counts, quiz pass
                             rate, average score, time spent \\
Temporal Signals       & 3 & certificate age, blockchain registration
                             delay, verification window duration \\
Verification Behaviour & 8 & total, successful, failed, and suspicious
                             attempt counts; unique IPs and user agents;
                             average response time; peak suspicion score \\
\bottomrule
\end{tabular}
\end{center}
\end{table}

\subsection{Fraud Scoring Pipeline}

The scoring pipeline is designed to be non-blocking: a failure in
the AI microservice never prevents certificate access. The sequence
is:

\begin{enumerate}
    \item On each certificate access, the current verification log
    count is compared against the cached log count stored on the
    \texttt{Certificate} document.
    \item If counts match, the cached fraud score is returned
    immediately without an AI service call.
    \item If counts differ, 24 features are extracted and sent to
    the AI microservice via HTTP POST with a 10-second timeout.
    \item A circuit breaker tracks consecutive failures. After two
    failures, the service is marked \textit{down} and a fallback
    heuristic score is returned for 60 seconds before retrying.
    \item On a successful AI response, the score, risk level, and
    SHAP signals are persisted back to the \texttt{Certificate}
    document for future cache hits.
\end{enumerate}

\subsection{SHAP Explainability}

SHAP TreeExplainer is applied to the trained XGBoost model to
decompose each fraud prediction into additive feature contributions.
For a certificate with predicted fraud log-odds $\hat{f}$:

\begin{equation}
\hat{f}(\mathbf{x}) = \phi_0 + \sum_{i=1}^{24} \phi_i(x_i)
\label{eq:shap}
\end{equation}

where $\phi_0$ is the base rate (mean model output over the training
set) and $\phi_i$ is the SHAP value for feature $i$, quantifying
its marginal contribution to the prediction relative to the base
rate. Features with $\phi_i > 0$ push the score toward fraud;
features with $\phi_i < 0$ push it toward normal. The five features
with the largest $|\phi_i|$ are returned to the admin dashboard
alongside their raw values and direction labels.

% ─────────────────────────────────────────────────────────────────
\section{Experimental Setup}

\subsection{Dataset Construction}

No publicly available dataset captures the specific fraud patterns
that arise in online certificate verification systems. A synthetic
dataset was therefore constructed using domain-informed generation
rules that reflect four empirically motivated fraud scenarios:

\textbf{Scenario A — Automated Verification Attacks:}
A certificate receives a high volume of verification requests
originating from a small number of IP addresses, with a large
fraction of requests failing and many flagged as suspicious by the
verification logger. This pattern is characteristic of credential
stuffing or bot-driven probing.

\textbf{Scenario B — Revoked Certificate Persistence:}
A certificate that has been administratively revoked continues to
receive verification attempts. The combination of revocation status
and ongoing verification activity is a strong fraud indicator.

\textbf{Scenario C — Blockchain Wallet Reuse:}
A single blockchain wallet address is associated with a large number
of certificates. Legitimate users rarely register more than one or
two certificates from the same wallet; a high reuse count suggests
bulk fraudulent issuance.

\textbf{Scenario D — Implausible Completion:}
A certificate is issued within minutes of course enrolment, with
near-zero time spent and suspiciously perfect quiz scores. This
pattern indicates automated course completion rather than genuine
learning.

The dataset contains 1,000 records: 850 normal and 150 fraudulent
(15\% fraud rate), split 80/20 for training and testing with
stratification. Class imbalance is addressed by setting XGBoost's
\texttt{scale\_pos\_weight} to $850/150 \approx 5.67$.

\subsection{Model Configurations}

\textbf{XGBoost:} 400 gradient-boosted trees, maximum depth 5,
learning rate 0.06, row subsampling 0.85, column subsampling 0.85,
minimum child weight 3, $\gamma = 0.1$, $\alpha = 0.1$,
$\lambda = 1.0$, objective \texttt{binary:logistic}.

\textbf{GraphSAGE GNN:} Two SAGEConv layers (hidden dimension 64,
then 32), dropout 0.3, 300 training epochs, Adam optimiser with
learning rate 0.005 and weight decay $10^{-4}$. Graph edges connect
certificates that share the same wallet reuse count bucket, both
have suspicious attempts, both have more failed than successful
verifications, and fall in the same unique-IP quartile. Early
stopping monitors validation F1 every 50 epochs.

\subsection{Evaluation Protocol}

Both models are evaluated on the same held-out 200-sample test set
using AUC-ROC, F1 score, precision, and recall. Five-fold
stratified cross-validation is additionally reported for XGBoost
to confirm generalisation across data partitions.

% ─────────────────────────────────────────────────────────────────
\section{Results and Analysis}

\subsection{Comparative Model Performance}

Table~\ref{tab:comparison} reports test-set metrics for both models.

\begin{table}[htbp]
\caption{Test-Set Performance: XGBoost vs.\ GraphSAGE GNN}
\label{tab:comparison}
\begin{center}
\begin{tabular}{lcccc}
\toprule
\textbf{Model} & \textbf{AUC} & \textbf{F1} & \textbf{Precision} & \textbf{Recall} \\
\midrule
XGBoost       & 1.0000 & 1.0000 & 1.0000 & 1.0000 \\
GraphSAGE GNN & 0.9967 & 0.9667 & 0.9667 & 0.9667 \\
\bottomrule
\end{tabular}
\end{center}
\end{table}

XGBoost achieves perfect separation on the test partition. GraphSAGE
attains AUC\,=\,0.997 and F1\,=\,0.967, indicating that 3.3\% of
fraudulent certificates are misclassified as normal. This recall gap
is attributable to sparse graph connectivity: with 1,000 nodes and
bucket-based edge construction, many certificates have only one or
two neighbours, limiting the GNN's ability to propagate fraud
signals across the graph. This finding is consistent with prior
work showing that tree-based models retain an advantage over GNNs
when graph structure is weak~\cite{ref9}.

\subsection{Confusion Matrix Analysis}

Table~\ref{tab:confusion} shows the XGBoost confusion matrix on the
200-sample test set (30 fraud, 170 normal).

\begin{table}[htbp]
\caption{XGBoost Confusion Matrix (Test Set, $n = 200$)}
\label{tab:confusion}
\begin{center}
\begin{tabular}{lcc}
\toprule
 & \textbf{Predicted Normal} & \textbf{Predicted Fraud} \\
\midrule
\textbf{Actual Normal} & 170 (TN) & 0 (FP) \\
\textbf{Actual Fraud}  & 0 (FN)   & 30 (TP) \\
\bottomrule
\end{tabular}
\end{center}
\end{table}

Zero false positives and zero false negatives on the test set
confirm that the four synthetic fraud patterns are well-separated
from normal behaviour in the 24-dimensional feature space. The
absence of false positives is particularly important in this
application domain, as incorrectly flagging a legitimate certificate
would damage user trust.

\subsection{SHAP Feature Importance}

Fig.~\ref{fig:shap} presents the mean absolute SHAP values across
all 1,000 certificates. The five most influential features are:

\begin{enumerate}
    \item \texttt{suspicious\_attempts} — the count of verification
    events flagged as suspicious by the request analyser. High
    values strongly increase the fraud score, directly capturing
    Scenario A behaviour.
    \item \texttt{max\_suspicious\_score} — the peak suspicion score
    assigned to any single verification event for this certificate.
    \item \texttt{failed\_verifications} — a high failure count
    relative to total attempts indicates systematic probing rather
    than legitimate employer verification.
    \item \texttt{is\_revoked} — a binary flag that, when set,
    produces a large positive SHAP contribution, capturing
    Scenario B directly.
    \item \texttt{wallet\_reuse\_count} — the number of certificates
    sharing the same blockchain wallet address, the primary signal
    for Scenario C.
\end{enumerate}

\begin{figure}[htbp]
\centerline{\includegraphics[width=\columnwidth]{shap_importance.png}}
\caption{Mean absolute SHAP values for the top 15 features. Bars
extending rightward indicate features that increase the fraud
probability; the ranking reflects average influence across all
1,000 certificates in the dataset.}
\label{fig:shap}
\end{figure}

\subsection{Cross-Validation Stability}

Five-fold stratified cross-validation on the full 1,000-sample
dataset yields AUC\,=\,1.000\,$\pm$\,0.000, confirming that the
model's performance is stable across all data partitions and is not
an artefact of a favourable train/test split.

\subsection{Operational Performance}

The caching mechanism reduces AI microservice calls by approximately
70\% in steady-state operation, since most certificate accesses
occur without new verification events. The circuit breaker ensures
that a cold-start delay on Render's free-tier infrastructure (which
spins down idle services after 15 minutes) does not propagate as
an error to end users — cached scores are served transparently
during the warm-up period.

% ─────────────────────────────────────────────────────────────────
\section{Discussion}

\subsection{Complementary Fraud Coverage}

Blockchain and machine learning address fundamentally non-overlapping
fraud vectors. The blockchain layer detects \textit{structural
fraud} — any modification to the certificate's immutable fields
(recipient identity, course, issue date) invalidates the on-chain
hash, providing cryptographic proof of tampering that requires no
ML inference. The ML layer detects \textit{behavioural fraud} —
patterns in how a certificate is used after issuance that indicate
misuse, even when the certificate itself is structurally intact.
A fraudster who obtains a legitimate certificate through academic
dishonesty would pass the hash check but would likely exhibit
anomalous verification behaviour detectable by the classifier.

\subsection{Synthetic Data Limitation}

The primary constraint of this study is that training data is
entirely synthetic. The four fraud scenarios are grounded in domain
knowledge but may not capture the full diversity of real-world
attack strategies. On real verification data, model performance is
expected to be lower than the reported AUC of 1.00, which reflects
the clean separability of synthetic patterns. The retraining
pipeline is specifically designed to address this: as the platform
accumulates real verification logs, the model can be retrained on
progressively more realistic data using the \texttt{retrain.py}
script.

\subsection{GNN Scalability Outlook}

Although XGBoost outperforms GraphSAGE at the current dataset scale,
the GNN's relative advantage is expected to increase as the platform
grows. At tens of thousands of certificates with denser graph
connectivity, the GNN's capacity to propagate fraud signals through
certificate clusters — for example, detecting a ring of certificates
all registered from the same wallet and verified from the same IP
subnet — may yield superior recall for coordinated fraud. This
represents a natural direction for future work.

\subsection{Deployment Constraints}

The system is deployed on Render's free tier, which imposes a
512\,MB memory limit per service. This constraint ruled out
\texttt{sentence-transformers} and \texttt{torch} as server-side
dependencies, as their combined footprint exceeds 2\,GB. The
circuit-breaker pattern was introduced specifically to handle the
15-minute spin-down behaviour of free-tier services gracefully,
ensuring that AI service unavailability never blocks certificate
issuance or verification.

% ─────────────────────────────────────────────────────────────────
\section{Conclusion}

This paper presented BlockShield, a production-deployed certificate
fraud detection system that integrates Ethereum blockchain
immutability with XGBoost-based behavioural scoring and SHAP
explainability. The system targets two distinct fraud vectors:
post-issuance data tampering, addressed through cryptographic hash
registration on the Sepolia testnet, and behavioural anomalies,
addressed through continuous ML scoring of 24 verification features.
Experimental evaluation on a 1,000-sample domain-informed dataset
demonstrates that XGBoost achieves AUC\,=\,1.00 and F1\,=\,1.00,
outperforming GraphSAGE GNN (AUC\,=\,0.997, F1\,=\,0.967) while
providing superior interpretability through per-certificate SHAP
explanations. The deployed system incorporates circuit-breaker fault
tolerance, score caching, and an automated retraining pipeline.
Future work will focus on retraining with real verification data,
evaluating GNN performance at larger graph scales, and extending
the feature set with temporal sequence features derived from
verification event streams.

% ─────────────────────────────────────────────────────────────────
\begin{thebibliography}{00}

\bibitem{ref1}
K. Smolenski and P. Barber, ``Blockcerts: An open infrastructure
for academic credentials on the blockchain,'' MIT Media Lab
Technical Report, 2016.

\bibitem{ref2}
A. Grech and A. F. Camilleri, ``Blockchain in education,''
Publications Office of the European Union, Luxembourg, 2017.

\bibitem{ref3}
T. Chen and C. Guestrin, ``XGBoost: A scalable tree boosting
system,'' in \textit{Proc. 22nd ACM SIGKDD Int. Conf. Knowledge
Discovery and Data Mining (KDD)}, San Francisco, CA, 2016,
pp.\ 785--794.

\bibitem{ref4}
J. Xia, Z. He, and W. Li, ``Credit card fraud detection using
gradient boosting and class imbalance handling,''
\textit{IEEE Access}, vol.\ 9, pp.\ 112054--112065, 2021.

\bibitem{ref5}
W. L. Hamilton, Z. Ying, and J. Leskovec, ``Inductive representation
learning on large graphs,'' in \textit{Advances in Neural
Information Processing Systems (NeurIPS)}, vol.\ 30, 2017,
pp.\ 1024--1034.

\bibitem{ref6}
S. M. Lundberg and S.-I. Lee, ``A unified approach to interpreting
model predictions,'' in \textit{Advances in Neural Information
Processing Systems (NeurIPS)}, vol.\ 30, 2017, pp.\ 4765--4774.

\bibitem{ref7}
A. Mahbooba, M. Timilsina, R. Sahal, and M. Serrano,
``Explainable artificial intelligence (XAI) to enhance trust
management in intrusion detection systems using decision tree
model,'' \textit{Complexity}, vol.\ 2021, Article ID 6634811, 2021.

\bibitem{ref8}
D. Bussmann, P. Nern, and C. Schwendner, ``Explainability for
machine learning models: From data transformations to model
predictions,'' \textit{Journal of Risk}, vol.\ 23, no.\ 6,
pp.\ 1--27, 2021.

\bibitem{ref9}
L. Grinsztajn, E. Oyallon, and G. Varoquaux, ``Why tree-based
models still outperform deep learning on tabular data,'' in
\textit{Advances in Neural Information Processing Systems
(NeurIPS)}, vol.\ 35, 2022, pp.\ 507--520.

\bibitem{ref10}
S. Nakamoto, ``Bitcoin: A peer-to-peer electronic cash system,''
\textit{Bitcoin.org}, 2008. [Online]. Available:
https://bitcoin.org/bitcoin.pdf

\bibitem{ref11}
N. Kshetri, ``Blockchain's roles in meeting key supply chain
management objectives,'' \textit{International Journal of
Information Management}, vol.\ 39, pp.\ 80--89, 2018.

\bibitem{ref12}
M. Sharples and J. Domingue, ``The blockchain and kudos: A
distributed system for educational record, reputation and reward,''
in \textit{Proc. 11th European Conf. Technology Enhanced Learning
(EC-TEL)}, Lyon, France, 2016, pp.\ 490--496.

\end{thebibliography}

\end{document}
